\documentclass[12pt, a4paper]{article}

% ── Packages ──────────────────────────────────────────────────────────────────
\usepackage[a4paper, margin=2.5cm]{geometry}
\usepackage{fontenc}
\usepackage{inputenc}
\setlength{\headheight}{15pt}
\usepackage{graphicx}
\usepackage{xcolor}
\usepackage{titlesec}
\usepackage{fancyhdr}
\usepackage{booktabs}
\usepackage{array}
\usepackage{tabularx}
\usepackage{enumitem}
\usepackage{parskip}
\usepackage{hyperref}
\usepackage{amsmath}
\usepackage{setspace}
\usepackage[numbers, compress]{natbib}
\usepackage{url}
\usepackage{caption}
\usepackage{subcaption}
\usepackage{float}
\usepackage{tcolorbox}

% ── Colours ───────────────────────────────────────────────────────────────────
\definecolor{usydnavy}{RGB}{0, 36, 77}
\definecolor{usydgold}{RGB}{228, 172, 43}
\definecolor{midgray}{RGB}{180, 180, 180}
\definecolor{placeholderbg}{RGB}{240, 244, 250}
\definecolor{placeholderborder}{RGB}{0, 36, 77}

% ── Figure path setup ─────────────────────────────────────────────────────────
% Image files live under ../images/test N - <description>/, so we add each test
% folder we draw figures from as a graphicspath entry. \includegraphics calls
% then refer to images by basename.
\graphicspath{%
  {../images/test 10 - current models rerun/}%
  {../images/test 11 - weather vs no wx t+120h/}%
  {../images/test 14 - drought onset 168h 24h tolerance/}%
  {../images/test 17 - station holdout/}%
  {../images/}%
}

% Figure helper. Usage: \realfig{filename.png}{label}{caption}
\newcommand{\realfig}[3]{%
  \begin{figure}[H]
    \centering
    \includegraphics[width=0.92\linewidth]{#1}
    \caption{#3}
    \label{fig:#2}
  \end{figure}%
}

% Retained for fallback if any image is missing during a build.
\newcommand{\figplaceholder}[3]{%
  \begin{figure}[H]
    \centering
    \begin{tcolorbox}[
      colback=placeholderbg,
      colframe=placeholderborder,
      boxrule=0.8pt,
      arc=3pt,
      left=6pt, right=6pt, top=6pt, bottom=6pt,
      width=\linewidth
    ]
      \centering
      \texttt{\textcolor{usydnavy}{[Figure placeholder: \textbf{#1}]}}\\[0.3em]
      \textcolor{midgray}{\small Replace with the rendered version of #1}
    \end{tcolorbox}
    \caption{#3}
    \label{fig:#2}
  \end{figure}%
}

% ── Hyperref setup ────────────────────────────────────────────────────────────
\hypersetup{
    colorlinks=true,
    linkcolor=usydnavy,
    urlcolor=usydnavy,
    citecolor=usydnavy,
    pdfauthor={Oscar Everett, James Robison, Angus Wyles, Byron Evans},
    pdftitle={ENGG2112 Final Report},
}

% ── Section formatting ────────────────────────────────────────────────────────
\titleformat{\section}
  {\large\bfseries\color{usydnavy}}
  {\thesection.}{0.6em}{}[\vspace{-0.3em}\textcolor{usydgold}{\rule{\linewidth}{1.5pt}}]

\titleformat{\subsection}
  {\normalsize\bfseries\color{usydnavy}}
  {\thesubsection.}{0.5em}{}

\titlespacing{\section}{0pt}{1.4em}{0.6em}
\titlespacing{\subsection}{0pt}{1em}{0.4em}

% ── Header / footer ───────────────────────────────────────────────────────────
\pagestyle{fancy}
\fancyhf{}
\fancyhead[L]{\small\textcolor{usydnavy}{ENGG2112: Final Report}}
\fancyhead[R]{\small\textcolor{usydnavy}{Semester 1, 2026}}
\fancyfoot[C]{\small\textcolor{midgray}{\thepage}}
\renewcommand{\headrulewidth}{0.4pt}
\renewcommand{\headrule}{\hbox to\headwidth{\color{usydgold}%
  \leaders\hrule height \headrulewidth\hfill}}

% ─────────────────────────────────────────────────────────────────────────────
\begin{document}

% ══════════════════════════════════════════════════════════════════════════════
%  COVER PAGE
% ══════════════════════════════════════════════════════════════════════════════
\begin{titlepage}
\begin{center}

  {\color{usydgold}\rule{\linewidth}{3pt}}
  \vspace{0.4cm}
  {\color{usydnavy}\rule{\linewidth}{1pt}}
  \vspace{1cm}

  \IfFileExists{USYD logo.png}{\includegraphics[height=2.8cm]{{USYD logo}.png}}{}

  {\Large\bfseries\color{usydnavy} The University of Sydney}\\[0.2em]
  {\normalsize\color{usydnavy} Faculty of Engineering}\\[0.1em]
  {\normalsize\color{midgray} ENGG2112: Multi-Disciplinary Engineering}

  \vspace{1.2cm}
  {\color{usydgold}\rule{0.6\linewidth}{1.5pt}}
  \vspace{0.8cm}

  {\LARGE\bfseries\color{usydnavy}
    SoilSync: AI-Powered Irrigation\\[0.3em] Decisions for Water-Scarce Farmland
  }

  \vspace{0.5cm}
  {\large\color{usydnavy} Final Report}

  \vspace{0.6cm}
  {\color{usydgold}\rule{0.6\linewidth}{1.5pt}}
  \vspace{1.0cm}

  {\large\bfseries\color{usydnavy} Team Members}\\[0.7em]
  \begin{tabular}{ll}
    \toprule
    \textbf{Name} & \textbf{Role} \\
    \midrule
    Oscar Everett  & ML Modelling \& Model Evaluation \\
    James Robison  & Project Lead \& Documentation \\
    Angus Wyles    & Domain Researcher \& Report Writing \\
    Byron Evans    & Data Engineer \& Analysis \\
    \bottomrule
  \end{tabular}

  \vspace{1.2cm}
  {\normalsize\color{usydnavy}
    Semester 1, 2026 \quad $\cdot$ \quad \today
  }

  \vfill

  {\color{usydnavy}\rule{\linewidth}{1pt}}
  \vspace{0.2cm}
  {\color{usydgold}\rule{\linewidth}{3pt}}

\end{center}
\end{titlepage}

% ══════════════════════════════════════════════════════════════════════════════
\section*{Executive Summary}
\addcontentsline{toc}{section}{Executive Summary}
% ══════════════════════════════════════════════════════════════════════════════

This report presents SoilSync, a machine learning system that forecasts agricultural
drought across 48 sensor stations in Kenya, Uganda and Rwanda. The work delivers
on the core objective stated in the A1 proposal, a comparison of an LSTM
network and a gradient-boosted regressor for short-range moisture forecasting,
and extends it along three axes. The architectural comparison was broadened
to four models (Random Forest, XGBoost, LSTM and a Temporal Fusion Transformer).
The problem framing was extended from pure regression to a dedicated drought
onset classification task that more closely mirrors the decision a farmer actually
faces. A cross-station holdout test was added to quantify how well the models
generalise to sensor locations they were not trained on.

Four models were trained on 1{,}103{,}461 quality-flagged hourly readings from
the International Soil Moisture Network, augmented by weather covariates from
the NASA POWER API. On the regression task at 72\,h, after upgrading the
sequence models to consume future weather through a dedicated encoder, the four
architectures converged to within $\Delta R^2 < 0.01$: LSTM achieved
$R^2 = 0.9367$, XGBoost $0.9302$, TFT $0.9272$ and Random Forest $0.9209$
(MAE between 0.0144 and 0.0159\,m$^3$/m$^3$). All four substantially exceed
the persistence baseline ($R^2 = 0.8822$), and the persistence improvement
widens monotonically with horizon.

The headline result is on the agriculturally critical task of \emph{drought
onset detection at seven days}: identifying timesteps where soil moisture is
currently adequate but will drop below 0.30\,m$^3$/m$^3$ within the next week.
Evaluated on 301 distinct onset events with a 24-hour success tolerance,
XGBoost catches 297 of 301 events (98.7\,\% event recall) at its best-F1
threshold, Random Forest catches 296 (98.3\,\%), LSTM catches 179 of 191 (93.7\,\%)
and TFT catches 167 of 191 (87.4\,\%). The persistence baseline catches zero
by construction; a linear-trend physical baseline catches 155 (51.5\,\%).
The estimated economic value of advance warning is \$58 to \$87 per hectare
per season in avoided yield losses for smallholder maize farmers, representing
16 to 24\,\% of gross seasonal revenue. On the station holdout test, Random Forest
and XGBoost generalise to 8 unseen stations with no measurable degradation
(holdout MAE within 0.001\,m$^3$/m$^3$ of the within-station benchmark). The
TFT failed catastrophically on this test until its static feature encoder was
regularised with dropout, an interpretable failure that we report in full. All
reported metrics use observed future weather as a proxy for a forecast and
therefore represent an upper bound on operational performance.

\newpage
\tableofcontents
\newpage

% ══════════════════════════════════════════════════════════════════════════════
\section{Project Overview, Objectives and Proposal Alignment}
% ══════════════════════════════════════════════════════════════════════════════

\subsection{Problem Statement and Users}

Smallholder farmers in Sub-Saharan Africa represent 75\,\% of those experiencing
famine and extreme poverty in the region, and in 95\,\% of cases depend entirely
on seasonal rainfall rather than managed irrigation \cite{worldbank2008}.
Agriculture accounts for 70\,\% of global freshwater withdrawals \cite{fao2017},
yet reactive irrigation systems trigger watering only after soil moisture drops
below a fixed threshold. This creates two structural inefficiencies: the lag
between a sensor alert and any moisture change can exceed 12 hours on farms with
slow-activating infrastructure, and the system has no awareness of incoming
rainfall, so fields may be irrigated even when significant precipitation is
forecast within 18 hours. Water spent irrigating ahead of forecast rainfall
cannot be recovered.

A multi-day moisture forecast addresses both deficiencies. It gives farmers
sufficient lead time to avoid redundant irrigation and to triage water use
across multiple plots when labour, pump capacity or water access is constrained.
The primary intended users are smallholder farmers and agricultural cooperatives
in East Africa, with secondary users including agronomic consultants, food
security NGOs and government water allocation agencies.

\subsection{Objectives}

The project pursued three objectives. The primary objective was to build and
compare machine learning models forecasting volumetric soil moisture at horizons
of 24 to 168 hours, and to determine which architecture generalises better
across geographic regions and seasons. The secondary objective was to assess
decision-support value through a dedicated drought \emph{onset} classifier:
correctly identifying upcoming moisture stress before it begins is of greater
agricultural consequence than reproducing the slow background trend that a
naive persistence model also captures. The third objective was to characterise
limitations honestly, particularly the substitution of observed future weather
for a real numerical weather prediction (NWP) forecast, which bounds reported
metrics as an upper limit rather than an operational estimate.

\subsection{Alignment with the A1 Proposal}

The A1 proposal committed to training and comparing an LSTM network and an
XGBoost regressor, evaluating both using RMSE and MAE across 24, 48 and 72-hour
horizons. The final project delivered on all of these commitments. Four extensions
were made beyond the original scope. First, a Random Forest and a Temporal Fusion
Transformer were added to the comparison, enriching the architectural analysis
without changing the original framing. Second, the evaluation horizon was extended
to 168 hours (seven days). Third, the sequence models (LSTM and TFT) were given
explicit access to future weather inputs through a dedicated encoder branch, so
the comparison against XGBoost is no longer asymmetric. Fourth, a standalone
drought onset classification task was added with three baselines (persistence,
linear-trend extrapolation, and the four ML models), reporting per-sample,
tolerant and event-level metrics. The geographic scope remained Uganda and Kenya
as proposed, with Rwanda added when the ISMN data pull returned stations there
as well. The data sources (ISMN and a weather API) and the evaluation metrics
(RMSE, MAE, recall, ROC AUC) matched the proposal exactly.

% ══════════════════════════════════════════════════════════════════════════════
\section{Data Sources, Preprocessing and Evaluation Strategy}
% ══════════════════════════════════════════════════════════════════════════════

\subsection{Soil Moisture Data}

Soil moisture data was sourced from the International Soil Moisture Network
(ISMN), the global reference database for in-situ soil moisture observations.
Two sub-networks were used: TAHMO (Trans-African Hydro-Meteorological Observatory)
contributing 46 stations and COSMOS (COsmic-ray Soil Moisture Observing System)
contributing 2 stations (KLEE and Mpala North). Geographic coverage spans Kenya
(approximately 30 stations), Uganda (approximately 10 stations) and Rwanda
(approximately 8 stations), with environments ranging from coastal lowlands
(Likoni, Base Titanium Dam) and highland forest (Kitabi College, IPRC Musanze)
to savanna (Maasai Mara cluster) and urban settings (Kampala KCCA school cluster).
Total temporal coverage spans 2016 to 2026 at hourly resolution. After quality
filtering, the dataset comprised 1{,}103{,}461 readings in volumetric water
content (m$^3$/m$^3$) at predominantly 0.2\,m sensor depth.

Raw data was distributed as ISMN \texttt{.stm} plain-text files and processed
through a bespoke cleaning pipeline. Only readings with a \texttt{G} (Good)
quality flag were retained. Sentinel no-data values of $-9999$ and $-999$ were
converted to \texttt{NaN}, duplicate timestamps were resolved by retaining the
first occurrence, and data was resampled to a uniform hourly grid by linear
interpolation capped at a maximum gap of six hours. Longer gaps remained as
\texttt{NaN} and were dropped at dataset construction time.

\subsection{Weather Data and Derived Features}

Weather covariates were obtained from the NASA POWER reanalysis API, which
provides complete global daily coverage with no missing values, making it
suitable for remote agricultural stations with limited ground-based
instrumentation. Seven variables were retrieved per station: mean, maximum and
minimum air temperature at 2\,m, relative humidity, corrected precipitation,
wind speed and all-sky surface shortwave solar radiation. A derived variable,
reference evapotranspiration (ET$_0$), was computed from these using the
Hargreaves-Samani equation:
\begin{equation}
  \text{ET}_0 = 0.0023 \times \frac{R_a}{0.408} \times (T_{\text{mean}} + 17.8)
  \times \sqrt{T_{\max} - T_{\min}}
\end{equation}
where $R_a$ is solar radiation converted to equivalent mm/day by dividing by
0.408. ET$_0$ quantifies atmospheric demand on soil water and is the primary
driver of moisture loss between rainfall events. Because weather data is daily
and soil moisture is hourly, daily weather values were broadcast across all
24 hours of each corresponding day. This approximation means each hourly timestep
within a day carries identical weather values, limiting the ability of sequence
models to exploit short-term weather dynamics.

\subsection{Train/Test Split and Evaluation Metrics}

The dataset was split per station using a chronological 80/20 partition. The
earliest 80\,\% of each station's readings formed the training set and the most
recent 20\,\% formed the test set. This per-station chronological split prevents
both temporal leakage (no future data is seen during training) and cross-station
contamination. Starting in Test~8, the same unified split was applied across
tabular and sequence builders, so XGBoost, Random Forest, LSTM and TFT are now
evaluated on the same per-station temporal partition rather than the global
80/20 index split used in earlier sequence runs.

Regression performance was evaluated using MAE (mean absolute error,
m$^3$/m$^3$), RMSE and $R^2$. The persistence baseline predicts that soil
moisture at time $t+h$ equals the soil moisture 24 hours prior, providing the
appropriate lower bound for practical utility.

For the drought onset task, three layers of metrics are reported:
\begin{itemize}
  \item \textbf{Per-sample exact}: each hourly test sample is classified as
        onset if current moisture is $\geq 0.30$\,m$^3$/m$^3$ and the moisture
        at exactly $t+h$ is $< 0.30$\,m$^3$/m$^3$. ROC AUC and PR AUC
        summarise discrimination across all probability thresholds.
  \item \textbf{Per-sample tolerant}: the sample is considered a true onset
        if drought occurs within a configurable window after the target time
        (24\,h in the final test). This avoids penalising the model for
        predicting an onset 18 hours before it materialises.
  \item \textbf{Event-level}: hourly samples are collapsed into contiguous
        warning/onset events. Event recall, event precision and event F1 are
        reported because a single missed event corresponds to one real
        agricultural failure, whereas a single missed hour does not.
\end{itemize}
Recall is prioritised over precision because a missed drought onset (false
negative) is agriculturally more costly than a false alarm.

The proposal defined success as MAE within 5\,\% of the observed moisture range.
The final XGBoost model achieved MAE of 0.0149\,m$^3$/m$^3$ at 72 hours, the
LSTM achieved 0.0144\,m$^3$/m$^3$. Across the test stations, the observed
moisture range spans approximately 0.05 to 0.55\,m$^3$/m$^3$, giving a range
of approximately 0.50\,m$^3$/m$^3$. The final MAE values therefore represent
roughly 2.9 to 3.0\,\% of this range, satisfying the success criterion with
margin.

% ══════════════════════════════════════════════════════════════════════════════
\section{Methodology and Technical Design}
% ══════════════════════════════════════════════════════════════════════════════

\subsection{Feature Engineering}

For the tree-based models, each row in the feature matrix represents one
timestep at one station. Soil moisture history is encoded as the current
reading plus five lag features (\texttt{sm\_lag\_24h} through
\texttt{sm\_lag\_120h}), covering a full five-day window. Rolling statistics
over 72 and 168 hours (mean and standard deviation, each shifted one hour to
prevent lookahead) encode trend and volatility. Temporal features (hour, month,
day of year) capture diurnal and seasonal cycles. Station metadata (latitude,
longitude, elevation, sensor depth) provides geographic and physical context.
The eight weather variables including ET$_0$ are appended as current-time
features.

A critical design decision was the inclusion of forward weather windows. To
simulate having a weather forecast, the model also receives actual observed
weather at each 24-hour step from $t+24$\,h up to the prediction horizon.
For the 72-hour horizon this adds $3 \times 8 = 24$ additional features; for
168\,h it adds $7 \times 8 = 56$. This constitutes an upper bound on
operational performance because real deployment would use NWP forecast values,
which carry their own prediction errors.

For sequence models (LSTM and TFT), a sliding window of 120 recent hourly
readings is passed as input. With weather channels, the input shape is
$(120, 9)$: 120 timesteps across nine channels comprising the soil moisture
value, seven NASA POWER variables and ET$_0$. Starting in Test~8, the LSTM
and TFT were extended with a dedicated future-weather encoder branch: the
$(h/24) \times 8$ matrix of future weather values is flattened, passed through
a small dense block (32-unit ReLU for LSTM; 16-unit ELU with layer
normalisation for TFT) and concatenated to the recurrent representation before
the output head. This closes the asymmetry that previously gave XGBoost an
unfair advantage over the sequence models. Each channel is normalised using
mean and standard deviation computed on the training split only.

\subsection{Model Architectures and Justification}

\textbf{Random Forest} uses an ensemble of 200 decision trees
(\texttt{max\_depth = 10}, \texttt{max\_features = sqrt}). As a bagging
ensemble it provides a strong and interpretable baseline for the tabular
feature set.

\textbf{XGBoost} uses gradient-boosted trees with histogram-based split finding.
Final parameters were \texttt{n\_estimators = 300}, \texttt{max\_depth = 5},
\texttt{learning\_rate = 0.05}, \texttt{subsample = 0.8} and
\texttt{colsample\_bytree = 0.8}. Unlike Random Forest, XGBoost builds trees
sequentially, each correcting the residuals of the previous ensemble.

\textbf{LSTM} was selected because recurrent architectures are well suited to
sequential time series where moisture at a given hour depends on prior
conditions. The final architecture is \texttt{LSTM(64, return\_sequences=True)}
$\to$ \texttt{Dropout(0.2)} $\to$ \texttt{LSTM(32)} $\to$ \texttt{Dropout(0.2)},
concatenated with the encoded future-weather vector, then
\texttt{Dense(16, ReLU)} $\to$ \texttt{Dense(1)}.

\textbf{TFT} extends the LSTM with a static feature encoder that injects
station metadata (latitude, longitude, elevation, depth) at every timestep,
and multi-head self-attention (4 heads, \texttt{key\_dim = 16}) that allows
the model to explicitly weight specific past timesteps. A gated residual
connection (sigmoid gate, add, layer normalisation) precedes the future-weather
fusion and output layers. Training uses a lower learning rate
($1\mathrm{e}{-4}$, with \texttt{clipnorm=1.0} and \texttt{ReduceLROnPlateau})
after earlier instability was traced to gradient explosion under the original
LSTM-equivalent settings.

\subsection{Drought Onset Classifier}

For the classification task a parallel pipeline reuses the same feature
builders, applies the same per-station temporal split, then \emph{filters}
both train and test sets to samples that are currently above the drought
threshold. The label is whether moisture at the target time crosses below
0.30\,m$^3$/m$^3$. Class imbalance is handled via
\texttt{class\_weight=`balanced'} for Random Forest and Keras models, and via
\texttt{scale\_pos\_weight} for XGBoost. Per-class thresholds are then swept
from 0.01 to 0.99 and three operating points are reported: the default
0.5 threshold, the best-F1 threshold, and the threshold that maximises recall
subject to a false-alarm-rate ceiling of 0.05. Two additional baselines were
added to give the classifier a meaningful comparator: a persistence baseline
(which by construction predicts no onset for currently-wet samples and therefore
achieves zero recall) and a \emph{linear-trend} physical baseline that
extrapolates the 24-hour slope to the forecast horizon and applies a sigmoid
to the crossing distance.

\subsection{Experiment Progression}

Fourteen experiments were run sequentially. Tests 2--7 (described in earlier
deliverables) established the regression pipeline, added weather features,
extended lag history, introduced the TFT, fixed normalisation leakage and
tuned hyperparameters, culminating in Test~7 with XGBoost at $R^2 = 0.9305$
at 72\,h. Test~8 unified the train/test split across all four models, giving
the sequence models the same per-station temporal partition. Test~8b added
oracle future-weather inputs to LSTM and TFT through the dedicated encoder
branch described above. Test~9 stabilised TFT training and added threshold
optimisation. Test~10 re-ran all four models under the unified pipeline and
added the first onset-event diagnostics. Test~11 ran a controlled
weather-vs-no-weather ablation at the 120-hour horizon across all four
architectures. Tests 12--14 then built out the dedicated drought onset
classifier: Test~12 introduced tolerant labels and event-level metrics, Test~13
added the linear-trend physical baseline, and Test~14 (the final configuration)
runs onset classification at the 168-hour horizon with a 24-hour success
tolerance, broadened training labels and the full six-classifier comparison.

% ══════════════════════════════════════════════════════════════════════════════
\section{Findings and Analysis}
% ══════════════════════════════════════════════════════════════════════════════

\subsection{Final Regression Performance at 72 Hours}

\begin{table}[H]
\centering
\caption{Regression performance at the 72-hour prediction horizon under the
unified per-station split, with future-weather inputs available to all four
models (Test~10). The four architectures converge to within
$\Delta R^2 < 0.02$. LSTM achieves the lowest MAE; XGBoost achieves the
highest $R^2$ on the tabular target while LSTM and TFT win on the sequence
target.}
\label{tab:final_results}
\begin{tabularx}{\linewidth}{lXXXX}
\toprule
\textbf{Model} & \textbf{MAE (m$^3$/m$^3$)} & \textbf{$R^2$} & \textbf{Recall} & \textbf{ROC AUC} \\
\midrule
Persistence    & 0.0187 & 0.8822 & \textendash & \textendash \\
Random Forest  & 0.0159 & 0.9209 & 0.977 & 0.970 \\
XGBoost        & 0.0149 & 0.9302 & 0.977 & 0.974 \\
LSTM           & \textbf{0.0144} & \textbf{0.9367} & 0.975 & \textbf{0.979} \\
TFT            & 0.0157 & 0.9272 & \textbf{0.978} & \textbf{0.980} \\
\bottomrule
\end{tabularx}
\end{table}

The predicted versus actual scatter plots in Figure~\ref{fig:predvsactual}
show that all four models now achieve tight clustering around the diagonal
across the full moisture range. The earlier finding that LSTM and TFT scattered
heavily at extreme low-moisture values no longer holds once the sequence
models receive future-weather inputs through the dedicated encoder branch.
This is corroborated by the residual distributions in
Figure~\ref{fig:residuals}: all four model residual distributions are tightly
centred near zero, and the four architectures are now meaningfully
indistinguishable on bulk regression metrics. The remaining differences emerge
on the drought-onset task discussed below.

\realfig{predicted_vs_actual.png}{predvsactual}{Predicted versus
actual soil moisture (m$^3$/m$^3$) for all four models at the 72-hour
horizon (Test~10). After unifying the temporal split and giving the sequence
models the same future-weather access as XGBoost, the four architectures all
cluster tightly around the diagonal.}

\realfig{residuals.png}{residuals}{Residual distributions for all four
models at the 72-hour horizon. All four are tightly centred near zero, with
LSTM showing the narrowest distribution.}

\subsection{Weather Feature Contribution (Test 11)}

A controlled weather-vs-no-weather ablation was run at the 120-hour horizon
across all four model families (Test~11). The results in
Table~\ref{tab:weather_ablation} show that adding weather and future-weather
inputs improves MAE by roughly 9 to 18\,\% across architectures, with XGBoost
gaining the most in absolute terms (MAE $0.0235 \to 0.0193$, $R^2$
$0.870 \to 0.902$). LSTM is the most weather-efficient model in relative
terms: it already performs strongly at 120\,h without weather ($R^2 = 0.891$)
and improves modestly to 0.901. Figure~\ref{fig:weatherimpact} visualises this
benefit as bar charts. The qualitative finding from the earlier report holds:
the weather contribution grows with horizon, because the autocorrelation of
the lag signal weakens while knowledge of future precipitation and
evapotranspiration becomes proportionally more informative.

\begin{table}[H]
\centering
\caption{Weather-vs-no-weather ablation at the 120-hour horizon across all
four models (Test~11). All four architectures improve substantially when given
forward weather information.}
\label{tab:weather_ablation}
\begin{tabularx}{\linewidth}{lXXXXX}
\toprule
\textbf{Model} & \textbf{Weather} & \textbf{MAE} & \textbf{RMSE} & \textbf{$R^2$} & \textbf{Recall} \\
\midrule
Random Forest & no  & 0.0236 & 0.0343 & 0.8650 & 0.970 \\
Random Forest & yes & 0.0208 & 0.0315 & 0.8856 & 0.977 \\
XGBoost       & no  & 0.0235 & 0.0336 & 0.8698 & 0.966 \\
XGBoost       & yes & 0.0193 & 0.0292 & 0.9021 & 0.972 \\
LSTM          & no  & 0.0214 & 0.0352 & 0.8914 & 0.984 \\
LSTM          & yes & 0.0208 & 0.0337 & 0.9009 & 0.974 \\
TFT           & no  & 0.0240 & 0.0390 & 0.8668 & 0.967 \\
TFT           & yes & 0.0203 & 0.0334 & 0.9021 & 0.962 \\
\bottomrule
\end{tabularx}
\end{table}

\realfig{weather_ablation_mae_bar.png}{weatherimpact}{Weather-vs-no-weather
ablation at the 120-hour horizon (Test~11). All four model families benefit
substantially from forward weather information, with XGBoost showing the
largest absolute MAE reduction.}

\subsection{Horizon Analysis}

\begin{table}[H]
\centering
\caption{XGBoost with weather versus the persistence baseline across all
prediction horizons (Test~10). The $R^2$ gap over persistence increases
monotonically with horizon length, confirming that the model adds
proportionally more value where naive forecasts are weakest.}
\label{tab:horizons}
\begin{tabularx}{\linewidth}{lXXXXXX}
\toprule
\textbf{Horizon} & \textbf{Persist MAE} & \textbf{XGB MAE} & \textbf{MAE Improv.} & \textbf{Persist $R^2$} & \textbf{XGB $R^2$} & \textbf{$R^2$ Gap} \\
\midrule
$t+24$\,h  & 0.0121 & 0.0079 & 34.7\,\% & 0.9363 & 0.9692 & $+$0.033 \\
$t+48$\,h  & 0.0157 & 0.0120 & 23.6\,\% & 0.9083 & 0.9476 & $+$0.039 \\
$t+72$\,h  & 0.0187 & 0.0149 & 20.3\,\% & 0.8822 & 0.9302 & $+$0.048 \\
$t+96$\,h  & 0.0212 & 0.0173 & 18.4\,\% & 0.8566 & 0.9159 & $+$0.059 \\
$t+120$\,h & 0.0235 & 0.0193 & 17.9\,\% & 0.8313 & 0.9032 & $+$0.072 \\
$t+168$\,h & 0.0273 & 0.0224 & 17.9\,\% & 0.7864 & 0.8813 & $+$0.095 \\
\bottomrule
\end{tabularx}
\end{table}

Figure~\ref{fig:horizoncomp} illustrates the degradation in MAE and $R^2$ with
horizon for both XGBoost and persistence. While absolute performance declines
for both, the gap between XGBoost and persistence widens monotonically: the
$R^2$ advantage grows from $+0.033$ at one day to $+0.095$ at seven days. The
high absolute $R^2$ values reflect the inherent autocorrelation of soil
moisture, a slow-moving physical variable. Persistence achieves $R^2 = 0.936$
at 24 hours with no modelling; the incremental improvement over this baseline,
not the absolute value, is the correct measure of model contribution.

\realfig{horizon_comparison.png}{horizoncomp}{MAE and $R^2$ versus
prediction horizon (24\,h to 168\,h) for XGBoost and the persistence baseline
(Test~10). The performance gap between XGBoost and persistence widens at longer
horizons.}

\subsection{Drought Onset Detection at Seven Days (Test 14)}

The decision-relevant task is not bulk regression but drought \emph{onset}
detection at long horizons: soil moisture is currently adequate but will drop
below 0.30\,m$^3$/m$^3$ within the forecast window. Test~14 evaluates this at
$t+168$\,h with a 24-hour success tolerance --- a true positive is recorded if
drought is observed at any point from 168\,h to 192\,h after the prediction is
made --- and collapses contiguous hourly predictions into events. This event
view is what matters operationally: a farmer cares about how many real onset
events the system catches, not how many individual hourly samples align.

\begin{table}[H]
\centering
\caption{Event-level drought onset detection at $t+168$\,h with 24\,h success
tolerance, evaluated at each model's best-F1 threshold (Test~14). Event recall
is the fraction of distinct real onset events the model issued at least one
warning for; event precision is the fraction of issued warning events that
correspond to a real onset. The tabular models share the test set; LSTM and
TFT share a slightly different test set due to the sequence builder's window
constraints.}
\label{tab:onset_events}
\begin{tabularx}{\linewidth}{lXXXXX}
\toprule
\textbf{Model} & \textbf{Actual events} & \textbf{Caught} & \textbf{Event recall} & \textbf{Event precision} & \textbf{Event F1} \\
\midrule
Persistence    & 301 & 0   & 0.000 & \textendash & 0.000 \\
Linear Trend   & 301 & 155 & 0.515 & 0.391 & 0.445 \\
Random Forest  & 301 & 296 & \textbf{0.983} & 0.778 & 0.869 \\
\textbf{XGBoost} & 301 & \textbf{297} & \textbf{0.987} & 0.747 & 0.850 \\
LSTM           & 191 & 179 & 0.937 & 0.722 & 0.816 \\
TFT            & 191 & 167 & 0.874 & 0.764 & 0.815 \\
\bottomrule
\end{tabularx}
\end{table}

The pattern is clear. Persistence is structurally unable to predict any onset
event because it forecasts no change from a currently-wet sample. A naive
physical baseline that extrapolates the recent moisture slope catches roughly
half the events but with poor precision (39\,\%, meaning over half its
warnings are false alarms). The four learned models all exceed 87\,\% event
recall. The tree-based models lead, with XGBoost catching 297 of 301 real
onset events (98.7\,\%) and Random Forest catching 296 (98.3\,\%). At their
best-F1 thresholds, both maintain roughly 75 to 78\,\% event precision ---
meaning that of every four issued warnings, three correspond to a genuine
upcoming onset.

At the more conservative operating point that bounds the false-alarm rate
below 0.05, all four learned models still detect roughly 90\,\% of real onset
events (Random Forest 92.0\,\%, XGBoost 92.0\,\%, LSTM 88.0\,\%, TFT 81.2\,\%),
demonstrating that high recall is not contingent on accepting an unmanageable
false-alarm load.

\realfig{drought_onset_curves.png}{onsetcurves}{Drought onset
classifier curves at $t+168$\,h with 24\,h tolerance (Test~14). Precision-recall
(left) and ROC (right) curves. XGBoost achieves the highest ROC AUC (0.946) and
PR AUC (0.819). The linear-trend baseline (gray) sits well below the learned
models, and persistence (black) lies on the no-skill diagonal.}

\realfig{drought_onset_precision_recall_bar.png}{onsetbar}{Onset
recall and precision at the best-F1 threshold for each classifier (Test~14).
The four learned models cluster in the same region. XGBoost and Random Forest
lead, the linear-trend baseline lags well behind, and persistence catches no
onset events.}

Figure~\ref{fig:recallhorizon} shows recall and ROC AUC trajectories from
earlier horizon sweeps. The recall advantage of the learned models over
persistence and the linear-trend baseline is consistent across all horizons,
confirming that the seven-day result is not a single-horizon artefact.

\realfig{recall_vs_horizon.png}{recallhorizon}{Onset recall and ROC
AUC versus prediction horizon (24\,h to 168\,h) for all models. The four
learned models maintain high recall across all horizons. The persistence
baseline is structurally unable to detect onsets in wet samples.}

\subsection{Feature Importance}

Figure~\ref{fig:featimportance} shows XGBoost feature importance for the
168-hour drought onset classifier (top 20 features from Test~14). The current
soil moisture reading dominates ($0.118$), followed by rolling moisture
statistics over 72 and 168 hours, then a cluster of future precipitation
features at $t+72$, $t+96$ and $t+120$ hours. This is intuitive: at the
seven-day horizon, the model needs to know how wet the soil is now, how its
recent trend has been moving, and how much rain is expected mid-window.
ET$_0$ features at multiple future horizons appear consistently in the top 20,
confirming that atmospheric evaporative demand is a real signal rather than
noise. Station metadata (latitude, longitude, elevation) contributes modest
but non-trivial importance, justifying its inclusion in the feature matrix.

\realfig{feature_importance.png}{featimportance}{XGBoost feature importance for
the 72-hour regression model (Test~10, top 20 features by gain). Current soil
moisture and rolling soil moisture statistics dominate, followed by future
precipitation and ET$_0$ features at the mid-window horizons. The same ranking
holds for the 168-hour onset classifier.}

\subsection{Cross-Station Generalisation (Test 17)}

The per-station temporal split used in Tests 8 to 14 trains and tests on the
same 48 sensors at different times. That is the right evaluation for whether
the model can forecast forward in time at known locations. It does not answer
the operationally important question of whether the model still works at new
locations the system was never trained on. Test 17 addresses this directly by
holding out 8 of the 48 stations from training entirely, stratified by country
(3 Kenya, 4 Uganda, 1 Rwanda), and evaluating all four architectures on the
held-out stations. The remaining 40 stations form the training set.

Table~\ref{tab:holdout} reports the headline numbers. Random Forest and XGBoost
generalise to unseen stations with no measurable degradation. Both achieve a
holdout MAE within 0.0006\,m$^3$/m$^3$ of their within-station MAE in Test~10,
and $R^2$ values within 0.02 of their within-station scores. The deployed
XGBoost classifier is therefore usable at new stations without retraining,
which is the most consequential finding for operational viability.

\begin{table}[H]
\centering
\caption{Regression performance on 8 held-out stations not seen during training
(Test~17, with weather features, $t+72$\,h). Within-station MAE is reproduced
from Table~\ref{tab:final_results} for comparison.}
\label{tab:holdout}
\begin{tabularx}{\linewidth}{lXXXX}
\toprule
\textbf{Model} & \textbf{Within-station MAE} & \textbf{Holdout MAE} & \textbf{Holdout $R^2$} & \textbf{Holdout ROC AUC} \\
\midrule
Random Forest & 0.0159 & 0.0154 & 0.943 & 0.987 \\
XGBoost & 0.0149 & 0.0145 & 0.946 & 0.987 \\
LSTM & 0.0144 & 0.0215 & 0.912 & 0.982 \\
TFT (with static-encoder dropout) & 0.0157 & 0.0264 & 0.872 & 0.968 \\
\bottomrule
\end{tabularx}
\end{table}

\realfig{mae_by_model.png}{holdoutbars}{Cross-station holdout MAE by model and
weather configuration (Test~17). Tree models are essentially unaffected by the
switch from within-station to cross-station evaluation. The LSTM bar is the
inverted case where adding weather features actually raises holdout MAE.}

The LSTM result depends strongly on the feature set. Without weather features,
LSTM achieves MAE 0.0139 and $R^2$ 0.942 on the 8 unseen stations, slightly
better than its within-station benchmark. Adding the weather and future-weather
inputs raises MAE to 0.0215 and lowers $R^2$ to 0.912. The likely explanation
is that the additional 24 weather channels per timestep give the network more
degrees of freedom to learn station-specific patterns that fail to transfer.
The same trade-off does not hurt the tree models, which still benefit from
weather (XGBoost MAE 0.0165 with no weather, 0.0145 with weather).

\subsubsection*{TFT failure mode and fix}

The TFT result required architectural investigation. In its original
configuration, the architecture used in Tests 8 through 14, the TFT scored MAE
0.0903 and $R^2 = -0.593$ with weather features on the holdout set. A negative
$R^2$ means the model performs worse than predicting the training mean for every
sample. Per-station holdout $R^2$ ranged from $-1.6$ to $-16.7$. ROC AUC for
the same predictions was 0.831, much higher than the negative $R^2$ would
suggest. This divergence between regression accuracy and ranking accuracy was
diagnostic. The predictions retained their relative ordering within each unseen
station but were systematically biased in absolute value.

The bias source was traced to the static feature encoder. The encoder consumes
(latitude, longitude, elevation, depth) through two dense layers and broadcasts
the result across every timestep of the input sequence. With no regularisation
it learns a per-station offset table, which works under the within-station
split but produces incorrect baselines when fed coordinates from stations it
never saw during training. Adding a single \texttt{Dropout(0.2)} layer inside
the static encoder forces the rest of the network to function with intermittently
dropped station identifiers during training, which breaks the offset-table
behaviour. With this one change the TFT achieves MAE 0.0264 and $R^2 = 0.872$
on the holdout set, which puts it back in the same performance band as the
other three models.

This finding is worth recording for two reasons. The first is that the same
architecture scored well in both Test~10 ($R^2 = 0.927$) and Test~14 (event
recall 0.874, event F1 0.815). Neither result gave any indication that the
model could not generalise to a new sensor. Within-station evaluation hides any
failure mode that depends on station identity, and classification thresholding
hides any failure mode that depends on absolute calibration. Both blind spots
applied to the original TFT evaluation. The second is that the fix is a single
dropout layer in a specific location, not a more capable architecture or more
data. The failure was localised to the static feature pathway, and so was its
remedy.
% ══════════════════════════════════════════════════════════════════════════════

\subsection{Architectural Comparison Revisited}

The original framing in earlier deliverables --- that XGBoost outperforms
sequence models because it has unfair access to explicit lag and future-weather
features --- no longer holds in the final pipeline. Once LSTM and TFT receive
future weather through a dedicated encoder branch (Test~8b onward) and once
the train/test split is unified across all four model families (Test~8), the
four architectures converge to within $\Delta R^2 < 0.02$ on the regression
task at 72 hours. LSTM in fact achieves the lowest MAE (0.0144\,m$^3$/m$^3$)
and the highest $R^2$ (0.9367), with TFT closely behind on ROC AUC (0.980).

Where the architectures \emph{do} differ is on the operational drought-onset
task. The tree-based models lead at both the per-sample level (XGBoost ROC AUC
0.946 vs.\ LSTM 0.895 and TFT 0.892 on tolerant labels) and the event level
(XGBoost event recall 0.987 vs.\ LSTM 0.937 and TFT 0.874 at best-F1). The
likely explanation is that the boosted trees can exploit thousands of explicit
lag and future-weather interactions through their tabular split structure,
whereas the sequence models compress the same information through a learned
recurrent representation that loses some of the sharp threshold-crossing
signals that determine onset. This is consistent with the broader literature
on tabular vs.\ sequence models for structured time-series forecasting.

The cross-station holdout result in Section~4.7 reinforces this picture and
adds nuance. Random Forest and XGBoost generalise to unseen stations almost
perfectly, with holdout MAE essentially identical to within-station MAE. The
LSTM generalises well without weather but degrades when weather is added,
which suggests the extra channels invite station-specific overfitting that the
tree models structurally avoid. The TFT in its original form did not generalise
at all, and only became competitive after the static encoder was regularised.
The hierarchy on the holdout test is therefore the same as on within-station
data (RF and XGBoost ahead of LSTM ahead of TFT), but the gap between tree
ensembles and the sequence models widens once generalisation to new sensors
is what is being measured.

\subsection{Limitations}

Six limitations bound the validity of the reported results. First, and most
critically, all experiments use observed future weather rather than a weather
forecast. In deployment, inputs at $t+24$ through $t+168$ hours would be drawn
from an NWP model carrying its own prediction errors, most significant at
longer horizons. The reported metrics are an upper bound, not an operational
estimate.

Second, the cross-station holdout test in Section~4.7 covers 8 stations from
the same two sensor networks (TAHMO and COSMOS) and the same three-country
region. Performance on a different sensor type, instrument calibration regime
or climate (for example South Asian rice paddies, North American corn belt,
semi-arid Sahel) is unknown and would require additional data to validate.
The result is that we can claim deployment readiness within the East African
TAHMO network but not, on the current evidence, globally.

Third, daily weather resolution broadcast to hourly soil moisture means each
LSTM and TFT timestep within a day carries identical weather values, limiting
exploitation of sub-daily dynamics.

Fourth, the 0.30\,m$^3$/m$^3$ drought threshold is a general guideline and is
not calibrated to specific crops or soil types at each station. Different
crops and soils have different wilting points; the threshold should ultimately
be parametrised per deployment.

Fifth, the model lacks soil texture data. Clay fraction and sand fraction
fundamentally determine how moisture moves through the soil profile; the
SoilGrids API provides this globally at no cost.

Sixth, all sensors operate at 0.2\,m depth. Root zone moisture integrated over
0 to 1\,m is more directly relevant to crop yield decisions than near-surface
readings alone.

\subsection{Stakeholders, Impact and Ethics}

The primary users are smallholder farmers and agricultural cooperatives
deciding when and how much to irrigate. Secondary stakeholders include
agronomic consultants, food security NGOs and government water allocation
agencies. A seven-day advance warning of moisture stress gives farmers time
to prioritise irrigation across multiple plots before crop stress begins,
which a threshold alarm system cannot provide. Importantly, the event-level
framing of the final classifier means the system issues a small number of
durable warnings (one per upcoming event) rather than thousands of correlated
hourly alerts.

Using conservative FAO-based assumptions for East African smallholder maize
production, the estimated economic value of the model over persistence is
\$58 to \$87 per hectare per season in avoided yield losses, representing
16 to 24\,\% of gross seasonal revenue on a \$360 per hectare base. This
estimate assumes the farmer has water access and acts on the warning; it does
not account for irrigation infrastructure costs or the cost of false alarms,
and should be treated as an illustrative order-of-magnitude figure.

The ethical obligations of deploying this system are substantive. ISMN data
skews toward North America and Europe; performance in Sub-Saharan Africa must
be explicitly validated and any degradation honestly reported. An incorrect
forecast could cost a farmer an entire season's crop, and silent failure could
cause systematic mismanagement across many farms before detection. Any
deployment must present prediction intervals and explicit operating points
(default, best-F1, low-false-alarm) rather than a single point estimate, so
that forecast uncertainty is communicated alongside the forecast itself. The
open data sources (ISMN, NASA POWER) and reproducible pipeline keep the
system accessible to users who cannot afford proprietary platforms. All data
used is publicly available and fully anonymised.

\subsection{Demo Application}

A live Plotly Dash application was built to demonstrate the final pipeline
end-to-end. The app loads the pre-trained XGBoost onset classifier and exposes
the seven-day drought outlook for each station in the network, with the
leaderboard, station drivers, network map and motion layer reorganised to
foreground the drought-onset narrative rather than raw moisture values. The
app illustrates how the operating-point selection (default, best-F1, low
false-alarm) would surface to a farmer-facing interface.

\subsection{Recommendations and Future Work}

The highest-priority extension is coupling the system with a real NWP forecast
model to replace observed future weather, converting the reported upper bound
into an operational performance estimate. The second priority is a broader
cross-region holdout that extends Section~4.7 beyond the TAHMO network and the
three current countries, ideally to a sensor network with different soil
profiles and climate. Adding soil texture features from SoilGrids, incorporating
monthly NDVI as a proxy for plant transpiration demand, implementing quantile
regression in XGBoost to produce prediction intervals on the regression task,
and disaggregating performance by season (long rains March to May, dry June
to August, short rains October to December) would each improve agricultural
specificity. For the sequence models, a richer fusion mechanism between the
recurrent representation and the future-weather encoder (cross-attention
rather than late concatenation) is a natural next step and could close the
remaining onset-recall gap to the tree-based models. The static encoder
diagnostic from Section~4.7 also suggests that any future model with static
covariates should be regularised in that pathway by default and stress-tested
under a cross-entity holdout before being trusted at new locations.

Commercially, the system occupies a position between free general-purpose
weather services and expensive proprietary precision agriculture platforms.
The open pipeline and freely accessible data sources make the approach
reproducible and extensible. A viable deployment pathway would integrate the
seven-day onset classifier with a simple mobile interface, delivering
event-level drought warnings and irrigation priority rankings directly to
farmers. Revenue could be structured as a cooperative subscription, with costs
shared across member farms.

% ══════════════════════════════════════════════════════════════════════════════
\newpage
\bibliographystyle{IEEEtran}
\bibliography{references}

% ══════════════════════════════════════════════════════════════════════════════
\newpage
\appendix
\section*{Appendix A: Team Organisation}
\addcontentsline{toc}{section}{Appendix A: Team Organisation}

Roles were divided based on each member's strengths and prior coursework,
agreed during an initial planning meeting.

\begin{table}[H]
\centering
\caption{Team roles and primary responsibilities.}
\begin{tabularx}{\linewidth}{llX}
\toprule
\textbf{Member} & \textbf{Role} & \textbf{Responsibilities} \\
\midrule
James Robison  & Project Lead      & Scoping, dataset sourcing, timeline management \\
Oscar Everett  & ML Engineer       & Model development, hyperparameter tuning, evaluation \\
Byron Evans    & Data Engineer     & Cleaning pipeline, feature engineering, dataset validation \\
Angus Wyles    & Domain Researcher & Literature review, domain context, report writing \\
\bottomrule
\end{tabularx}
\end{table}

Task tracking was managed through a shared Jira board and all code was version
controlled in a shared GitHub repository. The team met weekly to review
progress against the work plan proposed in A1.

\section*{Appendix B: Generative AI Usage}
\addcontentsline{toc}{section}{Appendix B: Generative AI Usage}

Generative AI tools were used in a supporting capacity during this project.
Claude (Anthropic) was used to assist with drafting and structuring report
sections, generating the \LaTeX{} template, and reviewing written content for
clarity and consistency. All technical results, model architectures,
experimental findings and quantitative claims in the body of the report are
drawn directly from the team's own experimental analysis. The team reviewed
and verified all AI-assisted drafts for accuracy before inclusion.

\end{document}
